\documentclass[12pt,a4paper]{article}
\usepackage[utf8]{inputenc}
\usepackage[french]{babel}
\usepackage[T1]{fontenc}
\usepackage{geometry}
\geometry{top=2.5cm, bottom=2.5cm, left=2.5cm, right=2.5cm}
\usepackage{graphicx}
\usepackage{hyperref}
\usepackage{booktabs}
\usepackage{xcolor}
\usepackage{titlesec}
\usepackage{fancyhdr}
\usepackage{lmodern}

% Configuration des couleurs
\definecolor{primary}{HTML}{0D6EFD}
\definecolor{secondary}{HTML}{495057}
\definecolor{lightbg}{HTML}{F8F9FA}

% Mise en page
\pagestyle{fancy}
\fancyhf{}
\rhead{\textcolor{secondary}{Méthodologie Scrum}}
\lhead{\textcolor{secondary}{EMSI 4IIR}}
\cfoot{\thepage}

\titleformat{\section}{\color{primary}\normalfont\Large\bfseries}{\thesection}{1em}{}
\titleformat{\subsection}{\color{secondary}\normalfont\large\bfseries}{\thesubsection}{1em}{}

\title{
    \vspace{2cm}
    \textbf{\Huge Organisation Agile du Projet}\\[0.5cm]
    \Large Plateforme de Gestion des Absences Multi-Rôles (JEE 8 / JPA)\\[1cm]
    \large\textbf{Rapport Scrum \& Gestion de Projet}
    \vspace{1cm}
}
\author{
    \textbf{Réalisé par :}\\
    Équipe de Projet — 4IIR\\[0.5cm]
    \textbf{Encadrant :}\\
    Pr. Houssam BAZZA
}
\date{Juin 2026}

\begin{document}

\maketitle
\thispagestyle{empty}
\newpage

\section{Introduction et Rôles Scrum}
Dans le cadre du développement de notre plateforme de gestion des absences, nous avons adopté la méthodologie agile \textbf{Scrum}. Ce cadre de travail nous a permis de structurer nos développements en cycles courts, de réagir rapidement aux contraintes techniques et de garantir un incrément logiciel de haute qualité.

\subsection{Composition de l'Équipe Scrum}
L'équipe est composée des rôles suivants :
\begin{itemize}
    \item \textbf{Product Owner (PO) :} Représente les besoins de l'administration scolaire et du corps enseignant (EMSI). Il est responsable de la spécification et de la priorisation des exigences (User Stories).
    \item \textbf{Scrum Master (SM) :} \textbf{Nizar Belkhadir} — Assure la bonne application des pratiques Scrum, anime les rituels agiles, résout les blocages techniques (Git, Docker, etc.) et maintient la vélocité de l'équipe.
    \item \textbf{Équipe de Développement (Development Team) :} \textbf{Nizar, Adam B., Adam A., Chafiq}. Équipe pluridisciplinaire responsable de la conception UML, du développement des composants front-end (JSP, CSS), back-end (Servlets, DAOs) et du déploiement.
\end{itemize}

\newpage
\section{Planification des Sprints}
Le projet a été développé sur une durée globale de trois semaines, découpée en trois Sprints successifs d'une semaine chacun.

\subsection{Sprint 1 : Conception, Persistance et Charte Graphique}
\begin{itemize}
    \item \textbf{Objectif :} Établir le modèle relationnel et poser les fondations de l'application.
    \item \textbf{Fonctionnalités et Tâches :}
    \begin{itemize}
        \item Analyse du besoin et conception des diagrammes UML (Cas d'utilisation, Classes, Séquence).
        \item Initialisation du projet Maven et configuration du gestionnaire de versions Git.
        \item Déploiement de la base de données PostgreSQL 15 sous conteneur Docker.
        \item Écriture des entités JPA persistantes (\texttt{Etudiant}, \texttt{Absence}, \texttt{Seance}, etc.).
        \item Développement de l'infrastructure de persistance sécurisée (\texttt{JPAUtil} et \texttt{ThreadLocal}).
        \item Création du Design System global en CSS (style Glassmorphism moderne).
    \end{itemize}
    \item \textbf{Incrément logiciel :} Schéma de base de données auto-généré et architecture de persistance opérationnelle.
\end{itemize}

\subsection{Sprint 2 : Cœur Applicatif et Sécurité des Accès}
\begin{itemize}
    \item \textbf{Objectif :} Implémenter l'authentification et les fonctionnalités principales d'appel et de consultation.
    \item \textbf{Fonctionnalités et Tâches :}
    \begin{itemize}
        \item Développement de la Servlet unique de routage (\texttt{AbsenceServlet}) appliquant le patron Front Controller.
        \item Gestion des sessions utilisateurs (\texttt{HttpSession}) et protection des URLs par rôle.
        \item Implémentation du module Enseignant : sélection de séance et feuille d'appel interactive avec enregistrement groupé.
        \item Implémentation du module Administrateur : registre des absences et CRUD manuel.
        \item Développement du moteur d'exportation de données au format CSV avec injection du BOM UTF-8.
    \end{itemize}
    \item \textbf{Incrément logiciel :} Application fonctionnelle permettant aux enseignants d'enregistrer les absences et aux administrateurs de les exporter.
\end{itemize}

\subsection{Sprint 3 : Workflows de Modération et Polissage}
\begin{itemize}
    \item \textbf{Objectif :} Finaliser les flux utilisateurs secondaires et réaliser la recette fonctionnelle.
    \item \textbf{Fonctionnalités et Tâches :}
    \begin{itemize}
        \item Implémentation du module Étudiant : statistiques d'assiduité et tableau de bord personnel.
        \item Workflow de justification : téléversement de fichiers binaires justificatifs (\texttt{MultipartConfig}) et modération administrative.
        \item Workflow de réclamation : soumission de contestations par l'étudiant et arbitrage administratif.
        \item Persistance locale de connexion avec cookies sécurisés « Remember Me ».
        \item Écriture du script de filtrage réactif Vanilla JS sur le client.
        \item Tests d'intégration et rédaction finale du rapport technique sous LaTeX.
    \end{itemize}
    \item \textbf{Incrément logiciel :} Produit fini stable, documenté et prêt pour la soutenance orale.
\end{itemize}

\newpage
\section{Rituels Scrum et Artefacts}
Afin de fluidifier les échanges et d'assurer une amélioration continue, l'équipe a tenu rigoureusement les cérémonies Scrum suivantes :

\subsection{Cérémonies Agile}
\begin{itemize}
    \item \textbf{Sprint Planning :} Réunion hebdomadaire d'une heure visant à définir l'objectif du sprint et à sélectionner les tâches du Product Backlog à accomplir.
    \item \textbf{Daily Stand-up :} Point quotidien de 15 minutes durant lequel chaque membre détaille ses actions de la veille, ses objectifs du jour et ses éventuels blocages.
    \item \textbf{Sprint Review :} Démo de fin de sprint devant le Product Owner pour valider la conformité de l'incrément développé.
    \item \textbf{Sprint Retrospective :} Réunion de fin de sprint pour identifier les points positifs et les axes d'amélioration méthodologique ou technique.
\end{itemize}

\subsection{Backlog de Produit et Exemples de User Stories}
\begin{table}[htbp]
    \centering
    \caption{Exemples clés de User Stories du Backlog}
    \vspace{0.3cm}
    \begin{tabular}{lp{6cm}p{5cm}}
        \toprule
        \textbf{Rôle} & \textbf{User Story} & \textbf{Objectif Métier} \\
        \midrule
        Professeur & En tant qu'enseignant, je veux faire l'appel sur une grille interactive & Enregistrer les absences rapidement. \\
        Étudiant & En tant qu'étudiant, je veux téléverser un certificat médical PDF & Justifier mes absences auprès de l'administration. \\
        Administrateur & En tant qu'administrateur, je veux exporter les absences en CSV & Analyser et éditer les données sous Microsoft Excel. \\
        Étudiant & En tant qu'étudiant, je veux contester une absence indue & Permettre sa suppression par l'administrateur. \\
        \bottomrule
    \end{tabular}
\end{table}

\end{document}
